\section{Discussion}

\subsection{Mitigations}
\label{sub:mitigations}

First and foremost, the vulnerability in Erlang SSH can be mitigated by disallowing all messages unrelated to the handshake phase before the completion of the first key exchange. This vulnerability is a straightforward layer violation, and thus easy to mitigate by safeguarding the implementation against unexpected message IDs.

\paragraph{Avoiding Strict KEX Violations}

Our analysis, as well as the subsequent responsible disclosure process, indicates that there is no single solution for preventing strict KEX violations due to fundamental differences in the code base of different implementations and the multitude of potential pitfalls. Instead, thorough testing of an implementation is required to rule out the presence of such violations. We observed that most implementers encountered significant challenges in verifying their patches because the implementation could not establish connections that violated strict KEX for negative testing. To resolve this issue, implementers may use our artifacts to verify the correctness of their strict KEX implementations.

\paragraph{Full Transcript Verification}

We observe that violations of the handshake hardening measures of strict KEX go unnoticed because they are not cryptographically enforced. Incorporating a full transcript verification, similar to the \texttt{verify\_data} approach of TLS 1.3~\cite{rfc8446}, would be one direction to harden the SSH handshake further against attacks like Terrapin. Bäumer~et~al. already discussed the applicability of a full transcript hash to mitigate the Terrapin vulnerability~\cite{USENIX:BauBriSch24}. Full transcript verification is an alternative approach to mitigating the Terrapin vulnerability, intended to supersede strict KEX, and does not prevent strict KEX violations. Our results emphasize the benefits of a full transcript hash; however, we also acknowledge the challenges involved in applying this idea.

Implementing a full transcript verification first requires a canonicalized message order since, in SSH, some messages do not have a well-defined order and may be sent simultaneously. For example, some servers may wait until they receive the client's \msgkexinit message, while others send their \msgkexinit immediately after the banner exchange. Similar to strict KEX, the client and server could then negotiate the use of full transcript verification in the \msgkexinit messages as pseudo-algorithm identifiers. The use of an extension as part of the \msgextinfo message is not possible, as this message is only exchanged after the handshake is concluded. If the extension is negotiated, the client extends its service request to include the value $MAC(K_x, H_c || H_s)$ where $K_x$ is a separate MAC key derived during key derivation with a different label for the sole purpose of authenticating the handshake. $H_c$ and $H_s$ are the unidirectional transcript hashes sent by the client and server, respectively, beginning at the very start of the connection up to and including the parties' \msgnewkeys message and computed with the hash function of the negotiated key exchange. Upon receiving the client's \msgservicerequest containing the transcript MAC, the server verifies it and aborts the connection if verification fails. Similarly, but using a different MAC key $K_y$, the server authenticates the transcript to the client in its \msgserviceaccept message, ensuring that both client and server acknowledged the transcript.

\paragraph{Limitations of Transcript Verification}

Note that the transcript verification approach does not ensure the correct rejection of invalid messages. For example, consider the Erlang vulnerability described in \autoref{sec:eval:sub:erlang}, which is triggered by an attacker sending an invalid channel open and channel request message. In this scenario, the attacker assumes the role of a client, giving them control over any session transcripts they issue. Neither including nor omitting these messages from the transcript on the server side would prevent the attack, as the attacker would always be able to produce a matching transcript of their own.
